% =============================================================================
% Anexa 2 — Formular de inscriere a temei lucrarii (TEMA LUCRARII DE DISERTATIE)
% Document oficial reprodus ca atare; se insereaza intre pagina de titlu si
% declaratia de onestitate. Pentru forma finala oficiala se inlocuieste cu
% scanul semnat in original.
% =============================================================================
\thispagestyle{empty}
{\setlength{\parindent}{0pt}\setlength{\parskip}{5pt}\footnotesize

% --- antet ---
\noindent
\begin{minipage}[t]{0.7\textwidth}
Universitatea Națională de Știință și Tehnologie POLITEHNICA București\\
Facultatea de Electronică, Telecomunicații și Tehnologia Informației\\
Programul de masterat \textbf{IISC}
\end{minipage}\hfill
\begin{minipage}[t]{0.25\textwidth}
\raggedleft\textbf{Anexa 2}
\end{minipage}

\vspace{1em}

\begin{center}
\textbf{TEMA LUCRĂRII DE DISERTAȚIE}\\[0.3em]
a masterandului \textbf{GEANĂ B.D. Tudor-Bogdan, 421-IISC}
\end{center}

\vspace{0.3em}

\textbf{1. Titlul temei:} Sisteme e-medicine de gestionare și monitorizare a
datelor epidemiologice

\textbf{2. Descrierea temei:}\\
Lucrarea prezintă proiectarea și realizarea unui sistem e-medicine pentru
analizarea și monitorizarea datelor epidemiologice, având la bază modelul
matematic SIR, implementat și testat comparativ în limbajele de programare R,
Julia, SAS și Python. Implementările sunt evaluate pe mai multe configurații
hardware, urmărindu-se comportamentul algoritmilor din punct de vedere al
performanței și al scalabilității. Datele epidemiologice utilizate sunt stocate
într-o bază de date distribuită YugabyteDB, rulată într-un cluster Kubernetes,
care asigură disponibilitate ridicată și distribuție eficientă a datelor. Pentru
testarea, compararea și analiza rezultatelor este dezvoltată o aplicație web
dedicată, compusă dintr-un API backend implementat în limbajul Go și o interfață
web realizată în Vue.js, prin intermediul căreia pot fi consultate rezultatele
rulării algoritmilor, seturile de date utilizate și graficele asociate acestora,
oferind o perspectivă unificată asupra performanței și comportamentului modelelor
epidemiologice analizate.

\textbf{3. Contribuția originală:}\\
Contribuția studentului constă în studiul și formalizarea modelului epidemiologic
SIR, urmate de implementarea acestuia în limbajele de programare R, Julia, SAS și
Python, într-o formă unitară care permite compararea directă a rezultatelor.
Studentul a proiectat și realizat scenarii de testare pe mai multe configurații
hardware, colectând și analizând metrici de performanță relevante. De asemenea,
studentul a proiectat arhitectura de stocare a datelor epidemiologice utilizând
YugabyteDB, implementată într-un cluster Kubernetes, asigurând distribuția și
disponibilitatea datelor. O contribuție majoră este dezvoltarea aplicației web
pentru testarea și compararea algoritmilor, incluzând implementarea unui API
backend în limbajul Go și a unei interfețe web în Vue.js, care permite
consultarea rezultatelor rulării algoritmilor, a seturilor de date utilizate și a
reprezentărilor grafice asociate, precum și integrarea tuturor componentelor
într-un sistem funcțional coerent.

\textbf{4. Resurse și bibliografie:}\\
Bibliografie selectivă\\
1. W. O. Kermack, A. G. McKendrick, A Contribution to the Mathematical Theory of Epidemics, Proc. R. Soc. A.\\
2. H. W. Hethcote, The Mathematics of Infectious Diseases, SIAM Review.\\
3. M. J. Keeling, P. Rohani, Modeling Infectious Diseases in Humans and Animals, Princeton University Press.\\
4. PostgreSQL Documentation: \texttt{https://www.postgresql.org/docs/}\\
5. YugabyteDB Documentation: \texttt{https://docs.yugabyte.com/}\\
6. Kubernetes Documentation: \texttt{https://kubernetes.io/docs/}\\
7. Go Documentation: \texttt{https://go.dev/doc/}\\
8. Vue.js Documentation: \texttt{https://vuejs.org/guide/}

\textbf{5. Data înregistrării temei:} 2026-01-22 12:26:55

\vspace{1.6em}

\noindent
\begin{tabular}{@{}p{0.52\textwidth}>{\raggedleft\arraybackslash}p{0.44\textwidth}@{}}
\textbf{Conducător(i) lucrare,}        & \textbf{Student,}            \\
Prof. dr. ing. Radu RĂDESCU            & GEANĂ B.D. Tudor-Bogdan      \\[2em]
\textbf{Responsabil program master,}   & \textbf{Decan,}              \\
Prof. dr. ing. Radu DOGARU             & Prof. dr. ing. Mihnea UDREA  \\
\end{tabular}

\vspace{1.6em}

\noindent Cod Validare: \textbf{e205b4fce0}

}
\clearpage
